\subsection{A single extracted registry across all protocol epochs}\label{sec:global:registry}
We now combine the per-block extractor into one
chronological procedure. The registry is an
\emph{analysis object}: it never appears in the
honest public protocol, and the extractor does
not open any commitment to the verifier.

\begin{definition}[Tagged bit-commitment registry]\label{def:global:registry}
A registered hidden field element is represented
by $n$ tagged Pedersen bit commitments
$(C_{x,0},\ldots,C_{x,n-1})$ and their
prescribed bitness OR blocks. Each tag uniquely
identifies a group element in the commitment
record; subsequent affine rows must reference
the \emph{same} tag and group element, not a
new value chosen under the same name.
A canonical extracted entry is
\[
 (C_{x,i};\,b_{x,i},r_{x,i}),\qquad
 C_{x,i}=g^{b_{x,i}}h^{r_{x,i}},
 \quad b_{x,i}\in\{0,1\}.
\]
The reconstructed field value is
$x^{\rm ext}=\sum_i\beta_i b_{x,i}$.
\end{definition}

\begin{definition}[Global checkpoint extractor]\label{def:global:extractor}
The extractor $\mathcal E_R$ follows exactly
the \emph{main} execution of the prescribed
protocol in Definition~\ref{def:global:epochs}.
Whenever a mandatory OR block accepts,
it immediately applies
Theorem~\ref{thm:global:one-block}
with at most $R$ side attempts, then restores
the main state. For a successful bitness
extraction it records the branch bit and
blinding in the registry.
For a successful aggregate-parity extraction
it records a permitted integer candidate
and its aggregate blinding, and verifies
its consistency with the already extracted
constituent bit entries.
If any repeated commitment tag has a different
extracted bit opening, or any aggregate
opening is inconsistent with the true
integer sum of its extracted constituent
bits, it outputs the derived value
$\log_g h$ using the reduction of
Chapter~10. Otherwise it continues
the main path without changing its
messages or challenge distribution.
The output is an accepting main transcript
with a complete consistent extracted
registry, a setup discrete logarithm,
or an extraction failure symbol.
\end{definition}

\begin{lemma}[A group element has a unique blinding for a fixed bit]\label{lem:global:unique-r}
If $\mathsf{Com}(b;r)=\mathsf{Com}(b;r')$,
then $r=r'$ in $\mathbb Z_P$.
\end{lemma}
\begin{proof}
Cancel the common factor $g^b$.
The equality becomes $h^r=h^{r'}$.
Because $h$ has prime order $P$,
this implies $r=r'\pmod P$.
\end{proof}

\begin{lemma}[Conflicting extracted bits reveal the setup logarithm]\label{lem:global:tag-conflict}
If two valid extracted openings
$(b,r)$ and $(b',r')$ of the same
tagged commitment have $b\neq b'$,
then $\mathcal E_R$ computes $\log_g h$.
\end{lemma}
\begin{proof}
The openings satisfy
$g^b h^r=g^{b'}h^{r'}$.
Lemma~\ref{lem:realize:binding-reduction}
provides the explicit modular inverse
formula for $\log_g h$.
\end{proof}

\begin{lemma}[False aggregate parity reveals the setup logarithm]\label{lem:global:row-conflict}
Suppose an accepted aggregate parity
block has extracted opening $(t,R_T)$,
where $t$ lies in the public permitted
candidate set, and all of its tagged
constituent bit commitments have
canonical extracted openings
$(b_i,r_i)$.
If the parity of $\sum_i b_i$
differs from the public parity target,
then $\mathcal E_R$ computes $\log_g h$.
\end{lemma}
\begin{proof}
Let $S=\sum_i b_i$ and
$R_{\rm true}=\sum_i r_i$.
Lemma~\ref{lem:realize:aggregate}
gives the valid aggregate opening
$(S,R_{\rm true})$. The OR fork gives
the second valid aggregate opening
$(t,R_T)$. Their parities differ,
so $S\neq t$ as ordinary integers.
The setup condition $P>|T|$
makes them distinct modulo $P$.
Theorem~\ref{thm:realize:parity-binding}
computes the forbidden logarithm.
\end{proof}

\begin{lemma}[Correct parity of every row gives all field-affine equations]\label{lem:global:all-affine}
Suppose all registered bit openings
are consistent, every prescribed affine
parity row is accepted and extracted,
and no aggregate row yields a setup
logarithm under
Lemma~\ref{lem:global:row-conflict}.
Then each prescribed field-affine
equation is true for the single
field tuple reconstructed from
the canonical bit registry.
\end{lemma}
\begin{proof}
Every accepted parity OR block
extracts an integer candidate
$t$ with the required target
parity. If the true integer sum
of its constituent canonical bits
had the opposite parity, then
Lemma~\ref{lem:global:row-conflict}
would return the setup logarithm,
contrary to assumption.
Thus every compiled row holds
over $\F_2$ for the canonical bits.
The exact affine-to-parity
equivalence
Theorem~\ref{thm:realize:affine-to-parity}
then makes every original
field-affine equation true for
the uniquely reconstructed
field values.
\end{proof}

\begin{theorem}[Global chronological extraction]\label{thm:global:chronological}
Assume the scheduling and checkpoint
model of Definitions~\ref{def:global:epochs}
and~\ref{def:global:access}; every
hidden bit appears under one stable
commitment tag, all required bitness
blocks are completed before the outer
challenge that uses that bit, and
all later compiled parity rows refer
to those existing tags.
For a main protocol with at most $J$
mandatory OR blocks, $\mathcal E_R$
has the following properties:
\begin{enumerate}[label=\textup{(\roman*)}]
\item The main visible transcript has
\emph{exactly} the law of the unmodified
honest-verifier interaction.
\item With probability at least
$1-\varepsilon_{\rm fork}$ on
the event of main acceptance,
understood in the unconditioned
joint-event sense
\[
 \Pr[\mathrm{main\ accepts}\ \wedge\
 \mathrm{any\ block\ extraction\ fails}]
 \le\varepsilon_{\rm fork},
\]
every mandatory accepted block yields
a valid opening.
\item Outside that failure event, either
$\mathcal E_R$ computes the setup
discrete logarithm, or the extracted
registry is consistent and every
accepted compiled field-affine
equation holds on one fixed field
tuple.
\item The field variables registered
in Epoch 0 are extracted
\emph{before} the actual sampling
of $\theta$; those registered in
Epoch 1 are extracted
\emph{before} the actual sampling
of $z$. The fresh outer challenges
retain their prescribed independent
uniform distributions.
\item The worst-case replay count is
at most $J(R+1)$ OR block executions.
\end{enumerate}
Here
$\varepsilon_{\rm fork}=
\min\{1,J(2^{-\kappa}+1/(R+1))\}$.
\end{theorem}
\begin{proof}
Apply the per-block procedure
of Theorem~\ref{thm:global:one-block}
at each accepted mandatory OR block.
Lemma~\ref{lem:global:main-law}
proves (i), because every side
replay is discarded and the
accepted main post-block state
is restored. The union bound
Corollary~\ref{cor:global:J-blocks}
proves the joint-event estimate
(ii). On its complement,
every accepted block has an
extracted permitted opening.
If two openings for the same
tag disagree on their bits,
Lemma~\ref{lem:global:tag-conflict}
computes the setup logarithm;
otherwise Lemma~\ref{lem:global:unique-r}
makes every tagged opening
unambiguous. If any extracted
aggregate candidate disagrees
with the bits of the same
registered commitments,
Lemma~\ref{lem:global:row-conflict}
computes the logarithm.
If neither happens,
Lemma~\ref{lem:global:all-affine}
establishes all prescribed
field-affine equations on
the canonical tuple. This
proves (iii).

For (iv), the protocol's Epoch 0
schedule finishes every required
bitness block and its immediate
side replay before the first
outer sampling step. Therefore
all Epoch 0 canonical entries are
already fixed when $\theta$
is drawn. The identical argument
applies to the Epoch 1 bitness
blocks before $z$.
Lemma~\ref{lem:global:outer-uniform}
shows that the outer challenges
remain independent uniforms.
Finally Lemma~\ref{lem:global:runtime}
gives (v).
\end{proof}

\begin{corollary}[No separately assumed global extractor for this affine compiler]\label{cor:global:discharged}
Within the \emph{classical checkpoint-replay
model} and the explicit registration schedule
of Theorem~\ref{thm:global:chronological},
the abstract global extraction obligation
of Definition~\ref{def:mv:global-coupling}
is realized by $\mathcal E_R$ with
the stated finite failure bound,
up to a discrete-log-breaking event.
This conclusion applies to the
bit-commitment/affine-proof component,
not to an unmodeled external KBFold
authenticated oracle or to quantum provers.
\end{corollary}
\begin{proof}
The five required properties
are precisely the output,
chronological consistency,
challenge independence and
same-main-path assertions of
Theorem~\ref{thm:global:chronological}.
The only exceptional events
are a failed side extraction
or a conflicting pair of
valid openings, which yields
the setup logarithm.
The oracle exclusion is part
of the checkpoint access model
and the declared transcript
domain, so it cannot be
eliminated by this argument.
\end{proof}

\begin{remark}[Conditioning is not being hidden]
The extraction error bound controls
the \emph{joint event} that an
accepted main run lacks a witness.
It does not claim a negligible
conditional failure probability
given acceptance if the adversary's
acceptance probability is itself
tiny. This distinction is needed
for the final soundness reduction.
\end{remark}
